\chapter{The Hybrid Collision-Avoidance Architecture}
\label{ch:ch24}
% Function names for pseudocode are prefixed with "Hybrid" to stay unique
% across the book.
\SetKwFunction{HybridLoop}{ControlCycle}
\SetKwFunction{HybridTrigger}{InsideHorizon}
\SetKwFunction{HybridIntended}{IntendedPositions}
\SetKwFunction{HybridReconnect}{Reconnect}
\SetKwFunction{HybridSegment}{SegmentFree}
\SetKwFunction{HybridReplan}{ReplanAndRecheck}
\SetKwFunction{HybridLocal}{LocalVelocity}
\SetKwFunction{HybridSearch}{SpaceTimeSearch}
\SetKwFunction{HybridDetect}{FirstConflict}
\SetKwFunction{HybridCbs}{LocalCBS}

\begin{objectives}
\begin{itemize}
  \item Draw the four-layer architecture as a block diagram with its data flows and rates, and state the inputs, outputs and guarantees of every layer.
  \item Describe the world model that the layers share, and convert a time-indexed grid path into a velocity reference and a continuous state back into a space-time start.
  \item Run the decision logic by hand: the state machine Nominal / Avoiding / Reconnecting / Replanning, the safety-horizon trigger built from a predicted trajectory and its uncertainty, the reconnection search, the replanning trigger and the conflict re-check.
  \item Say precisely which guarantees of \cbs, \orca and space-time \astar survive the integration, which do not, and which experiment of \cref{ch:ch25} measures the gap.
  \item Build the capstone prototype of the study plan (\cref{ch:appA}, Week~12) from the modules of the earlier chapters, and recognise the four integration pitfalls that break it.
\end{itemize}
\end{objectives}

% ---------------------------------------------------------------------
\section{Why this matters for a drone swarm}
\label{sec:ch24-motivation}

Every earlier chapter solved one problem well. \cbs and \ecbs
(\cref{ch:ch09,ch:ch10}) compute conflict-free, time-indexed routes for the
drones you control, but only for the world they were given before take-off.
The Kalman filter and the learned predictors (\cref{ch:ch18,ch:ch20}) tell you
where an unknown drone will be in the next few seconds, with an honest
covariance, but they do not decide anything. \orca and DWA
(\cref{ch:ch13,ch:ch14}) choose a safe velocity for the next tenth of a second,
but they have no memory, no plan and no notion of a swarm. \dstarlite and
\astar (\cref{ch:ch05,ch:ch04}) repair a route on a changed map, but a route
for one drone is worthless if it crosses the routes of the others. The
consensus and MPC controllers (\cref{ch:ch23,ch:ch21}) keep a formation and a
radio link, as long as nobody has to swerve.

The picture that motivates this chapter is the one of \cref{ch:ch01}: three
delivery drones fly a planned mission through a district, and a fourth drone
that nobody planned for crosses their lanes. The nominal plan must not be
thrown away at the first sign of trouble, because it is the only thing that
keeps the three from colliding with \emph{each other}; it must not be
followed blindly either, because it knows nothing about the intruder. The
\textbf{hybrid architecture}\index{hybrid architecture} of the training plan
(\cref{ch:appA}) resolves this with one rule: \emph{a global planner for the
routes, a local method for the surprises, a replanner only when the local
response is insufficient}. This chapter turns that sentence into an
executable design: a block diagram with rates, a shared world model, a
state machine, four precise decision rules with pseudocode, a fully worked
scenario run by the book's own code, and an honest account of what the
composed system does and does not guarantee.

The chapter is organised as follows. \Cref{sec:ch24-intuition} gives the
architecture and the interfaces, \cref{sec:ch24-world} the world model,
\cref{sec:ch24-logic} the state machine and the per-drone control loop.
\Cref{sec:ch24-horizon,sec:ch24-reconnect,sec:ch24-replan} define the three
decisions that drive the transitions: the safety-horizon trigger, the
reconnection, and the replanning with its conflict re-check.
\Cref{sec:ch24-constraints,sec:ch24-uncertainty} add the formation and
communication constraints and the treatment of prediction uncertainty.
\Cref{sec:ch24-scenario} runs the worked scenario, \cref{sec:ch24-guarantees}
states the guarantees, \cref{sec:ch24-implementation} describes the
implementation and its pitfalls, and \cref{sec:ch24-research} lists the
research directions of the training plan with concrete formulations. The
exercises end with the capstone build of Week~12.

% ---------------------------------------------------------------------
\section{The idea in plain words}
\label{sec:ch24-intuition}

Think of the four layers as four people with different time budgets
(\cref{fig:ch24-architecture}). The \emph{global planner} works before
take-off, with seconds or minutes, and produces one time-indexed path
$\pi_i$ per drone; it is the only component that reasons about all drones at
once, and its output is the contract the drones fly by. The \emph{prediction
layer} watches the sensors ten times a second and maintains, for each
unknown object, an estimate with a covariance and a short predicted
trajectory. The \emph{local safety layer} also runs ten times a second; it
receives a preferred velocity and the predicted obstacles and returns the
closest velocity that is safe for the next few seconds. The \emph{replanning
layer} runs only on demand: when a drone has been pushed so far off its
route that it cannot get back safely, it computes a new route from where the
drone actually is, on a map in which the predicted intruder positions are
blocked and the other drones' routes are reserved.

The glue between them is the \textbf{executive}\index{executive}: a small
state machine that runs in every drone at the control rate and decides which
layer is in charge. While no predicted conflict lies inside a short
\emph{safety horizon}, the executive follows the nominal path and the local
layer is idle. When the prediction says that the intruder will come too
close within the horizon, the local layer takes the velocity command, and the
nominal path is kept in memory. When the danger has passed, the executive
tries to \emph{reconnect}: it picks a waypoint a little further along the
nominal path and checks whether a straight flight to it is safe. Only if that
fails does it call the replanner, and after any replan it re-checks the new
route against the routes of the other drones, because a route that is safe
against the intruder may now cross a teammate.

\begin{figure}[tb]
  \centering
  \inputfigure{ch24/architecture}
  \caption{The four layers, the executive and the shared world model, with
  the data that flows between them and the rate at which each block runs.
  What to notice: the global planner and the replanner work on the
  grid-and-time side of the world model, the local layer and the predictor on
  the continuous side; the executive sits in the middle and is the only
  block that converts between the two. The dashed arrow is the conflict
  re-check that closes the loop back to the global planner.}
  \label{fig:ch24-architecture}
\end{figure}

\begin{keyidea}
Keep the global plan as the contract between the drones and treat every
deviation from it as temporary. A deviation is handled by the cheapest layer
that can handle it: first a velocity change for a few seconds, then a
reconnection to the plan, then a new plan for one drone, and only then a new
plan for several. Each escalation is triggered by a condition you can write
down, and each returns the drone to the contract or replaces the contract
with one that has been checked again.
\end{keyidea}

\Cref{tab:ch24-layers} is the architecture as a set of interfaces. The rates
are those of the code of this chapter and of a typical small quadrotor: a
control cycle of $\dt = 0.1$\,s, a prediction horizon of a few seconds, and a
grid time step of one second, which is the time a drone needs to cross one
cell at cruise speed.

\begin{table}[tb]
  \centering\footnotesize
  \caption{Responsibilities and interfaces of the layers. Each layer is a
  module with the inputs and outputs listed; the rate column is the one used
  in the worked scenario of \cref{sec:ch24-scenario}.}
  \label{tab:ch24-layers}
  \begin{tabularx}{\textwidth}{@{}lYYYl@{}}
  \toprule
  Layer (chapter) & Inputs & Outputs & Guarantee & Rate\\
  \midrule
  1 Global swarm planner: \cbs, \ecbs (\cref{ch:ch09,ch:ch10}) & inflated grid, starts, goals, swarm constraints & one time-indexed path $\pi_i$ per drone, $\sumcost$, $\makespan$ & conflict-free; optimal or $w$-suboptimal & once before take-off; on request after a failed local repair\\
  2 Prediction: Kalman filter, EKF/UKF/PF, LSTM (\cref{ch:ch18,ch:ch19,ch:ch20}) & noisy observations $\meas_k$ of every unknown object & estimate $\hat{\state}_k$, predicted means $\hat{\pos}_{T+h}$ and covariances $\mat{\Sigma}_h$ for $h\,\dt \le T_{\mathrm{pred}}$ & calibrated ellipse per step (if the model holds) & every control cycle, 10\,Hz\\
  3 Local safety: \orca, DWA (\cref{ch:ch13,ch:ch14}) & own state, teammates' states, predicted intruder disc, preferred velocity $\vel^{\mathrm{pref}}$ & velocity command $\vel^{\mathrm{new}}$, feasibility flag & pairwise safety for $\ttc$ under its assumptions & every control cycle, 10\,Hz\\
  4 Replanning: \dstarlite, \astar (\cref{ch:ch05,ch:ch04}) & current cell and start time, reservation table of the others, blocked cells per time layer & new path $\pi_i'$ from the current cell, or failure & conflict-free against the reserved paths & on demand, at most one per $T_{\mathrm{replan}}$\\
  Executive (this chapter) & everything above, formation offsets, $R_{\mathrm{comm}}$ & state, $\vel^{\mathrm{pref}}$, triggers, re-check & the composition of \cref{sec:ch24-guarantees} & every control cycle, 10\,Hz\\
  \bottomrule
  \end{tabularx}
\end{table}

% ---------------------------------------------------------------------
\section{The world model shared by the layers}
\label{sec:ch24-world}

The global and replanning layers live on a grid with discrete time; the
local and prediction layers live in continuous space and time. The two
descriptions must refer to the same world, the same clock and the same
frame, and the executive must convert between them in both directions.

\begin{definition}[Hybrid world model]
\label{def:ch24-world}
The \textbf{hybrid world model}\index{world model!hybrid}\index{hybrid
architecture!world model} consists of
\begin{enumerate}
  \item a grid $G = (V, E)$ of square cells of side $\ell$ (\cref{ch:ch02}),
  inflated by the drone radius, with a time step $\Delta T = \ell / v_{\mathrm{cruise}}$;
  the cell $(x, y)$ has its centre at $\vect{c}(x,y) = ((x + \tfrac12)\ell, (y + \tfrac12)\ell)$;
  \item for every drone $a_i$ a \textbf{time-indexed path}\index{time-indexed path} $\pi_i$ with a
  \textbf{start time}\index{path!start time $t_0$} $t_0^i \in \N$: $\pi_i[j]$ is the cell
  occupied at grid time $t_0^i + j$, and the drone parks at $\pi_i[T_i]$ from
  $t_0^i + T_i$ on (stay-at-target, \cref{ch:ch07});
  \item the \textbf{reservation table}\index{reservation table!with time offsets} $R$ of the
  other drones, holding the pairs (cell, grid time), the reversed moves and the
  parked goals of their paths in absolute time, exactly as in \cref{ch:ch04,ch:ch08};
  \item the \textbf{blocked layers}\index{blocked layer} $\mathcal{B}_t \subseteq V$: the cells that the
  predicted intruder may occupy at grid time $t$;
  \item the continuous states: own position $\pos_i(t)$ and velocity
  $\vel_i(t)$, the teammates' states, and for each intruder the estimate
  $\hat{\state}$ with predicted means $\hat{\pos}_{T+h}$ and covariances
  $\mat{\Sigma}_h$ at the look-ahead $s = h\,\dt$ for $0 \le s \le T_{\mathrm{pred}}$ (\cref{ch:ch20}).
\end{enumerate}
One clock $t$ (seconds since mission start) and one frame serve all of them;
grid time $t_{\mathrm{grid}} = t / \Delta T$.
\end{definition}

In the worked scenario $\ell = 1$\,m, $v_{\mathrm{cruise}} = 1$\,m/s, so
$\Delta T = 1$\,s and grid time equals time in seconds; the control cycle is
$\dt = 0.1$\,s.

\paragraph{From a path to a velocity reference.}
The flight controller needs a continuous reference. The
\textbf{reference position}\index{reference position} of drone $a_i$ at time
$t$ is the linear interpolation between the cells of its path,
\begin{equation}
  \pos^{\mathrm{ref}}_i(t) = \vect{c}(\pi_i[j]) + (s - j)\bigl(\vect{c}(\pi_i[j+1]) - \vect{c}(\pi_i[j])\bigr),
  \qquad s = \frac{t}{\Delta T} - t_0^i,\quad j = \lfloor s \rfloor,
  \label{eq:ch24-reference}
\end{equation}
clamped to the first cell for $s \le 0$ and to the last cell for $j \ge T_i$.
The \textbf{preferred velocity}\index{preferred velocity!from a path} that the
executive hands to the local layer is a pure pursuit of the reference a short
lookahead $t_{\mathrm{la}}$ ahead,
\begin{equation}
  \vel^{\mathrm{pref}}_i(t) = \frac{\pos^{\mathrm{ref}}_i(t + t_{\mathrm{la}}) - \pos_i(t)}{t_{\mathrm{la}}},
  \qquad \text{clipped to } \norm{\vel^{\mathrm{pref}}_i} \le v_{\mathrm{cruise}} \text{ while the path is unfinished},
  \label{eq:ch24-vpref}
\end{equation}
with $t_{\mathrm{la}} = 0.5$\,s in the code. A drone that sits on its reference
receives exactly the velocity of the reference; a drone that has been pushed
off receives a velocity that steers it back, and the pull grows with the
distance. The \textbf{drift}\index{drift from the plan}
$d_i(t) = \norm{\pos_i(t) - \pos^{\mathrm{ref}}_i(t)}$ is the distance from
the reference and is the quantity the executive monitors while the local
layer is in charge.

\paragraph{From a continuous state back to the grid.}
The replanner needs a space-time start. The executive uses the cell that
contains the drone, $c_0 = (\lfloor p_x/\ell \rfloor, \lfloor p_y/\ell \rfloor)$,
or its nearest free neighbour if that cell is blocked, and the \textbf{start
time}\index{replanning!start time $t_{\mathrm{start}}$}
$t_{\mathrm{start}} = \lceil t/\Delta T + \tfrac12 \rceil$, the next grid time
that leaves at least half a step to fly to the centre of $c_0$. The new path
is stored as $\pi_i' = (c_0, \pi_i'[1], \dots)$ with $t_0^i = t_{\mathrm{start}} - 1$,
so that \cref{eq:ch24-reference} interpolates from the current cell into the
plan. Everything a teammate reserves is expressed in absolute grid time
through its own $t_0$; the executive never stores a path without its start
time, which is the mistake of the third pitfall of
\cref{sec:ch24-implementation}.

\paragraph{The intruder on the grid.}
The prediction is a sequence of discs. At every grid time $t' \ge
t_{\mathrm{start}}$ inside the prediction horizon, the executive blocks the
cells whose centre lies within $R_{\mathrm{safe}} + \kappa_p\,\sigma(h) + \ell/2$
of the predicted mean, where $h$ is the prediction step that matches $t'$;
\cref{sec:ch24-uncertainty} defines the inflation. These are the layers
$\mathcal{B}_{t'}$; they change every cycle, which is the situation
\dstarlite was built for (\cref{ch:ch05}).

% ---------------------------------------------------------------------
\section{The decision logic: a state machine and a control loop}
\label{sec:ch24-logic}

The five steps of the training plan (follow, avoid, reconnect, replan,
re-check) become four states of the executive and the transitions of
\cref{fig:ch24-state-machine}. Each drone runs its own copy; the only shared
data are the paths in the world model.

\begin{figure}[tb]
  \centering
  \inputfigure{ch24/state-machine}
  \caption{The per-drone state machine. What to notice: every transition is
  labelled with a condition the executive can evaluate from the world model;
  the two self-loops are where the time is spent; and Reconnecting is the
  state through which every deviation returns to Nominal, whether the plan
  was kept, shifted or replaced.}
  \label{fig:ch24-state-machine}
\end{figure}

\begin{description}[style=nextline]
  \item[Nominal.] The drone follows $\pi_i$ through \cref{eq:ch24-vpref}
  with the formation and communication terms of \cref{sec:ch24-constraints}.
  The local layer is armed but sees only the teammates, and since the
  teammates follow a conflict-free plan its half-planes are inactive.
  \item[Avoiding.] The local layer commands the velocity. The intruder enters
  as a disc with full responsibility and an uncertainty-inflated radius, the
  teammates as reciprocal discs. The formation term is switched off; the
  nominal path is kept.
  \item[Reconnecting.] The drone flies a straight segment to a waypoint
  $\pi_i[k]$ of its plan, timed to arrive at grid time $t_0^i + k + \Delta$
  with a delay $\Delta \ge 0$; when it gets there it is Nominal again.
  \item[Replanning.] The replanning layer computes $\pi_i'$ from the current
  cell; the result is re-checked against the others. While no path exists the
  drone hovers under the protection of the local layer and retries.
\end{description}

\begin{table}[tb]
  \centering\small
  \caption{Transitions of the state machine, their conditions and the action
  taken. The parameters are listed in \cref{tab:ch24-parameters}.}
  \label{tab:ch24-transitions}
  \begin{tabularx}{\textwidth}{@{}llYY@{}}
  \toprule
  From & To & Condition & Action\\
  \midrule
  Nominal & Avoiding & a predicted conflict is inside the horizon: $t_c \le \tau_h$ (\cref{def:ch24-inside}) & remember $t_{\mathrm{enter}} = t$; keep $\pi_i$\\
  Avoiding & Reconnecting & no predicted conflict for $n_{\mathrm{clear}}$ consecutive cycles, and \HybridReconnect finds $(k, \Delta)$ & target $\pi_i[k]$ at $t_0^i + k + \Delta$; if $\Delta > 0$ shift the remainder of $\pi_i$ and re-check\\
  Avoiding & Replanning & drift $d_i > d_{\max}$, or $t - t_{\mathrm{enter}} > t_{\mathrm{avoid}}$, or the local program was infeasible, or \HybridReconnect fails & release the formation slot; keep hovering safely until a path exists\\
  Reconnecting & Nominal & the target waypoint time has been reached & clear the target\\
  Reconnecting & Avoiding & a predicted conflict is inside the horizon again & as Nominal $\to$ Avoiding\\
  Reconnecting & Replanning & the segment to the target is no longer free & as above\\
  Replanning & Reconnecting & \HybridReplan returned a path and the re-check passed (possibly after a local \cbs) & target the first cell of $\pi_i'$ at $t_{\mathrm{start}}$\\
  Nominal & Reconnecting & another drone's re-check changed $\pi_i$ (partial replanning) & target the first cell of the new path\\
  \bottomrule
  \end{tabularx}
\end{table}

\Cref{alg:ch24-loop} is the control loop that implements the table; it is
called once per cycle for every drone with the freshest prediction. The
functions it calls are the subject of the next three sections:
\HybridTrigger is \cref{alg:ch24-trigger}, \HybridReconnect is
\cref{alg:ch24-reconnect}, \HybridReplan is \cref{alg:ch24-replan}, and
\HybridLocal is the \orca program of \cref{ch:ch13} (or the DWA search of
\cref{ch:ch14}) with the half-planes chosen by the argument list.

\begin{algorithm}[htb]
\caption{The per-drone control cycle of the executive}
\label{alg:ch24-loop}
\KwIn{drone $a_i$ with state, path $\pi_i$, start time $t_0^i$; time $t$; prediction $(\hat{\pos}_{T+h}, \mat{\Sigma}_h)_h$ or \KwNil; states of the teammates}
\KwOut{the velocity command for this cycle}
\Fn{\HybridLoop{$a_i$, $t$}}{
  $(\mathit{inside}, t_c) \gets$ \HybridTrigger{$a_i$, $t$, prediction}\label{alg:ch24-loop:trigger}\;
  \If{state $=$ Nominal \KwAnd $\mathit{inside}$}{state $\gets$ Avoiding; $t_{\mathrm{enter}} \gets t$; $n_{\mathrm{ok}} \gets 0$\label{alg:ch24-loop:enter}\;}
  \If{state $=$ Avoiding}{
    $n_{\mathrm{ok}} \gets 0$ \textbf{if} $\mathit{inside}$ \textbf{else} $n_{\mathrm{ok}} + 1$\;
    \If{$d_i(t) > d_{\max}$ \KwOr $t - t_{\mathrm{enter}} > t_{\mathrm{avoid}}$ \KwOr the last local program was infeasible}{state $\gets$ Replanning\label{alg:ch24-loop:bound}\;}
    \ElseIf{$n_{\mathrm{ok}} \ge n_{\mathrm{clear}}$}{
      $(k, \Delta) \gets$ \HybridReconnect{$a_i$, $t$, prediction}\label{alg:ch24-loop:reconnect}\;
      \lIf{$(k,\Delta) =$ failure}{state $\gets$ Replanning}
      \Else{state $\gets$ Reconnecting; target $\gets (\vect{c}(\pi_i[k]),\, t_0^i + k + \Delta)$\;
        \lIf{$\Delta > 0$}{$\pi_i \gets \pi_i[k{:}]$; $t_0^i \gets t_0^i + k + \Delta$; re-check as in \cref{alg:ch24-replan}\label{alg:ch24-loop:shift}}}
    }
    \lElse{\KwRet \HybridLocal{$\vel^{\mathrm{pref}}_i$ without formation term; teammates reciprocal; intruder with full share, inflated radius at $t_c$}\label{alg:ch24-loop:avoid}}
  }
  \If{state $=$ Replanning}{
    \lIf{$t - t_{\mathrm{replan}} \ge T_{\mathrm{replan}}$ \KwAnd \HybridReplan{$a_i$, $t$, prediction} succeeds}{$t_{\mathrm{replan}} \gets t$; state $\gets$ Reconnecting; target $\gets (\vect{c}(c_0),\, t_{\mathrm{start}})$\label{alg:ch24-loop:replan}}
    \lElse{\KwRet \HybridLocal{$\vect{0}$; teammates and intruder as in Avoiding}\tcp*[f]{hover safely, retry next cycle}\label{alg:ch24-loop:hover}}
  }
  \If{state $=$ Reconnecting}{
    \lIf{$\mathit{inside}$}{state $\gets$ Avoiding; \KwRet the velocity of line~\ref{alg:ch24-loop:avoid}\label{alg:ch24-loop:again}}
    \lIf{$t \ge$ target time}{state $\gets$ Nominal; clear the target\label{alg:ch24-loop:back}}
  }
  \KwRet \HybridLocal{$\vel^{\mathrm{pref}}_i$ with formation and communication terms; teammates reciprocal; no intruder}\label{alg:ch24-loop:nominal}\;
}
\end{algorithm}

Three details of the loop deserve a comment. First, the trigger on
line~\ref{alg:ch24-loop:trigger} is evaluated in every state, so that a
reconnecting drone that flies into a new prediction returns to Avoiding on
line~\ref{alg:ch24-loop:again} before it is hit. Second, the Avoiding state
has two exits and the bounded one, line~\ref{alg:ch24-loop:bound}, is the
important one: without a bound on the drift and on the time, the local layer
can carry the drone arbitrarily far from its plan, and the reservation table
of the others, which assumes the plan, becomes fiction (first pitfall of
\cref{sec:ch24-implementation}). Third, the velocities returned on
lines~\ref{alg:ch24-loop:avoid}, \ref{alg:ch24-loop:hover} and
\ref{alg:ch24-loop:nominal} all pass through the local layer; even the
nominal flight is filtered by the teammates' half-planes, which costs nothing
while the plan is conflict-free and protects the drone when a teammate has
deviated.

\begin{table}[tb]
  \centering\small
  \caption{Parameters of the executive, their meaning, their value in the
  worked scenario and their effect. Distances in metres, times in seconds.}
  \label{tab:ch24-parameters}
  \begin{tabularx}{\textwidth}{@{}llY@{}}
  \toprule
  Parameter & Value & Meaning and effect\\
  \midrule
  $\dt$; $\Delta T$ & $0.1$; $1$ & control cycle; grid time step $= \ell/v_{\mathrm{cruise}}$\\
  $r_A$, $r_B$, $\Delta r$ & $0.3$, $0.3$, $0.1$ & drone radius, intruder radius, safety margin; $R_{\mathrm{safe}} = r_A + r_B + \Delta r = 0.7$\\
  $\tau_h$ & $3$ & safety horizon $=$ \orca horizon $\ttc$; larger $=$ earlier, gentler, more false alarms\\
  $\kappa_p$ & $2.45$ & inflation factor of the $95\,\%$ disc; larger $=$ safer and more conservative\\
  $T_{\mathrm{pred}}$ & $5$ & prediction horizon; must cover $\tau_h$ and the replanning layers\\
  $n_{\mathrm{warm}}$ & $10$ & filter updates before a prediction is trusted (\cref{ch:ch18})\\
  $n_{\mathrm{clear}}$ & $5$ & clear cycles before reconnection is attempted; hysteresis against chattering\\
  $K$; $\Delta_{\max}$ & $6$; $1$ & waypoints searched forward; largest admissible delay in grid steps\\
  $d_{\max}$; $t_{\mathrm{avoid}}$ & $3.5$; $8$ & drift and time bounds of Avoiding; the escalation to Replanning\\
  $T_{\mathrm{replan}}$ & $1$ & shortest interval between two replans of one drone\\
  $R_{\mathrm{comm}}$; $\ell_{\mathrm{act}}$ & $8.5$; $7.5$ & communication range; link length at which the range term acts\\
  $k_F$ & $0.5$ & formation correction gain of \cref{eq:ch24-formation}\\
  \bottomrule
  \end{tabularx}
\end{table}

% ---------------------------------------------------------------------
\section{The safety horizon and the time-to-collision trigger}
\label{sec:ch24-horizon}

The first decision is whether an intruder is a problem \emph{now}. It is
answered by comparing where the drone intends to be with where the intruder
is predicted to be, over a short window, with the prediction uncertainty
taken into account.

\paragraph{The intended motion.}
Let $\tilde{\pos}_i(s)$ be the \textbf{intended position}\index{intended
position} of drone $a_i$ a look-ahead $s \ge 0$ from now: the point reached
by flying at cruise speed along the polyline that starts at $\pos_i(t)$,
passes through the reconnection target if there is one, and then follows the
remaining waypoints $\vect{c}(\pi_i[k_0]), \vect{c}(\pi_i[k_0+1]), \dots$ of
the plan, where $k_0 = \lceil t/\Delta T - t_0^i \rceil$ is the first waypoint
not yet passed. For a drone on its reference this is the reference itself;
for a drone that has deviated it is the path it would take home. The
prediction layer supplies $\hat{\pos}_{T+h}$ and $\mat{\Sigma}_h$ at the
same look-aheads $s = h\,\dt$ (\cref{ch:ch20}), and we write
$\sigma(h) = \sqrt{\lambda_{\max}(\mat{\Sigma}_h)}$ for the largest
standard deviation of the predicted position.

\begin{definition}[Predicted conflict inside the safety horizon]
\label{def:ch24-inside}
Let $R_{\mathrm{safe}} = r_A + r_B + \Delta r$ be the required separation
between a drone of radius $r_A$ and an intruder of radius $r_B$, with a
margin $\Delta r \ge 0$, and let $\kappa_p$ be the confidence factor of
\cref{sec:ch24-uncertainty}. The \textbf{margin}\index{margin function}
at look-ahead $s = h\,\dt$ is
\begin{equation}
  m_i(s) = \norm{\tilde{\pos}_i(s) - \hat{\pos}_{T+h}} - \kappa_p\,\sigma(h) - R_{\mathrm{safe}}.
  \label{eq:ch24-margin}
\end{equation}
A predicted conflict is \textbf{inside the safety horizon}\index{safety
horizon}\index{hybrid architecture!safety horizon} $\tau_h$ if
$m_i(s) < 0$ for some $s \in [0, \tau_h]$, and the \textbf{predicted time to
collision}\index{time to collision!predicted} is the first such look-ahead,
\begin{equation}
  t_c = \min\set{s \in [0, \tau_h] : m_i(s) < 0}.
  \label{eq:ch24-tc}
\end{equation}
In practice $s$ runs over the multiples of $\dt$ up to $\tau_h$.
\end{definition}

The definition subtracts the inflation from the separation rather than
adding it to the radius; the two are the same thing, and
\cref{sec:ch24-uncertainty} explains why $\kappa_p\,\sigma(h)$ is the right
amount. \Cref{alg:ch24-trigger} is the test, and \cref{fig:ch24-horizon}
shows it at the moment it fires in the worked scenario: the predicted
separation of drone~A falls to $1.53$\,m at $s = 3.0$\,s while the threshold
$R_{\mathrm{safe}} + \kappa_p\sigma$ has grown to $0.70 + 0.86 = 1.56$\,m; one
cycle earlier the margin at the same look-ahead was still $+0.01$\,m.

\begin{algorithm}[htb]
\caption{The safety-horizon trigger}
\label{alg:ch24-trigger}
\KwIn{drone $a_i$, time $t$, prediction $(\hat{\pos}_{T+h}, \mat{\Sigma}_h)$ for $h = 0, \dots, \lfloor T_{\mathrm{pred}}/\dt \rfloor$ or \KwNil}
\KwOut{$(\mathit{inside}, t_c)$}
\Fn{\HybridTrigger{$a_i$, $t$, prediction}}{
  \lIf{prediction $=$ \KwNil}{\KwRet $(\KwFalse, \KwNil)$\tcp*[f]{no intruder, or the filter is still warming up}}
  $H \gets \min(\lfloor \tau_h/\dt \rfloor,\ \lfloor T_{\mathrm{pred}}/\dt \rfloor)$\;
  $(\tilde{\pos}_i(0), \dots, \tilde{\pos}_i(H\dt)) \gets$ \HybridIntended{$a_i$, $t$}\label{alg:ch24-trigger:intended}\;
  \For{$h \gets 0$ \KwTo $H$}{
    $m \gets \norm{\tilde{\pos}_i(h\dt) - \hat{\pos}_{T+h}} - \kappa_p \sqrt{\lambda_{\max}(\mat{\Sigma}_h)} - R_{\mathrm{safe}}$\label{alg:ch24-trigger:margin}\;
    \lIf{$m < 0$}{\KwRet $(\KwTrue, h\dt)$\label{alg:ch24-trigger:fire}}
  }
  \KwRet $(\KwFalse, \KwNil)$\;
}
\end{algorithm}

\begin{figure}[tb]
  \centering
  \inputfigure{ch24/horizon}
  \caption{The trigger of drone~A in the worked scenario at $t = 3.5$\,s.
  The predicted separation between A's intended positions and the intruder's
  predicted mean (solid blue) decreases with the look-ahead; the threshold
  $R_{\mathrm{safe}} + \kappa_p\sigma(s)$ (orange) grows with it because the
  prediction becomes less certain. What to notice: the two curves meet
  exactly at the horizon $s = \tau_h = 3$\,s, where the margin is
  $-0.03$\,m; one control cycle earlier (dashed) the separation was $0.05$\,m
  larger and the trigger stayed silent. The trigger fires at the horizon,
  not at the collision: three seconds remain to act.}
  \label{fig:ch24-horizon}
\end{figure}

\paragraph{Choosing the horizon.}
Three quantities constrain $\tau_h$. First, the horizon $\ttc$ of the local
layer: \orca's half-plane of \cref{ch:ch13} forbids the relative velocities
that collide within $\ttc$, so the local layer only starts to steer when
$t_c \le \ttc$. With $\tau_h \ge \ttc$ control is handed over before the
local layer has to act, and its first correction is small; with
$\tau_h < \ttc$ the hand-over comes late and the first correction is a jump.
The code uses $\tau_h = \ttc$. Second, the dynamics: a drone with
acceleration limit $a_{\max}$ needs $v/a_{\max}$ seconds to stop from speed
$v$ and covers the stopping distance $d_{\mathrm{stop}}(v) = v^2/(2a_{\max})$
of \cref{ch:ch02} while doing so; changing from one velocity of the disc
$\norm{\vel} \le v_{\max}$ to any other takes at most $2v_{\max}/a_{\max}$.
The horizon must leave time for the local layer's command to be executed,
\begin{equation}
  \tau_h \;\ge\; \dt + \frac{v_{\max}}{a_{\max}}
  \quad\text{(a stop suffices)},\qquad
  \tau_h \;\ge\; \dt + \frac{2v_{\max}}{a_{\max}}
  \quad\text{(a reversal may be needed)}.
  \label{eq:ch24-tauh}
\end{equation}
Third, the prediction: $\tau_h \le T_{\mathrm{pred}}$, and
$\kappa_p\sigma(\tau_h)$ must stay well below the spacing of the lanes,
otherwise the trigger fires for intruders in the next lane. In the scenario
$v_{\max}/a_{\max} = 1.5$\,s, $\ttc = 3$\,s and $\kappa_p\sigma(3\,\mathrm{s})
= 0.86$\,m, so $\tau_h = 3$\,s satisfies the first two conditions and the
inflation is still below the lane spacing of $2$\,m. A last condition ties
the horizon to the sensor: the intruder must be detected and tracked before
it enters the horizon, which needs a detection range of at least
$(v_{\max} + v_B)(\tau_h + n_{\mathrm{warm}}\dt) + R_{\mathrm{safe}} + \kappa_p\sigma(\tau_h)$,
about $7.5$\,m in the scenario, where the intruder is first seen at $6$\,m by
drone~B and enters A's horizon $2.2$\,s later.

\begin{pitfall}[A trigger without hysteresis chatters]
The margin at the horizon changes by a few centimetres per cycle, and noise
in the estimate moves it back and forth across zero. If the executive left
Avoiding in the first cycle without a predicted conflict, it would re-enter
it in the next one, and the drone would alternate between the \orca velocity
and the pursuit of its plan ten times a second. The counter
$n_{\mathrm{clear}}$ in \cref{alg:ch24-loop} is the remedy: the drone leaves
Avoiding only after $n_{\mathrm{clear}}$ consecutive clear cycles. A second
hysteresis is built into the trigger itself, because entering uses the
inflated radius at $t_c$ while the local layer keeps the intruder's
half-plane for the whole of $\ttc$.
\end{pitfall}

% ---------------------------------------------------------------------
\section{Reconnecting to the nominal path}
\label{sec:ch24-reconnect}

When the conflict has cleared, the cheapest repair is no repair at all: fly
back to the plan and resume it. The question is \emph{where} to rejoin, and
the answer is a search along the plan.

\begin{definition}[Admissible reconnection]
\label{def:ch24-reconnection}
A \textbf{reconnection}\index{reconnection}\index{hybrid
architecture!reconnection} of drone $a_i$ at time $t$ is a pair $(k, \Delta)$
with $k \ge k_0$ a waypoint index and $\Delta \in \set{0, \dots, \Delta_{\max}}$
a delay in grid steps such that the straight flight from $\pos_i(t)$ to
$\vect{c}(\pi_i[k])$, arriving at $t_k = (t_0^i + k + \Delta)\Delta T$,
(i) needs a speed $\norm{\vect{c}(\pi_i[k]) - \pos_i(t)}/(t_k - t) \le v_{\max}$,
(ii) keeps a clearance of at least $r_A$ from the static obstacles,
(iii) keeps the margin of \cref{eq:ch24-margin} non-negative against the
prediction at every sampled time, and (iv) stays at least
$R_{\mathrm{mate}} = 2r_A + \Delta r$ from the intended positions of every
teammate at the same times. Its \textbf{extra cost} is $\Delta$: the
arrival time of $a_i$ increases by $\Delta$ grid steps, and $\sumcost$ and
$\makespan$ by at most that.
\end{definition}

\begin{algorithm}[htb]
\caption{Reconnection search along the nominal path}
\label{alg:ch24-reconnect}
\KwIn{drone $a_i$ at $\pos_i(t)$ with path $\pi_i$, start time $t_0^i$ and last index $T_i$; prediction; look-ahead $K$; largest delay $\Delta_{\max}$}
\KwOut{$(k, \Delta)$ or \emph{failure}}
\Fn{\HybridReconnect{$a_i$, $t$, prediction}}{
  $k_0 \gets \max(0, \lceil t/\Delta T - t_0^i \rceil)$\label{alg:ch24-reconnect:k0}\;
  \For{$\Delta \gets 0$ \KwTo $\Delta_{\max}$\tcp*[f]{prefer no delay}}{
    \For{$k \gets k_0$ \KwTo $\min(k_0 + K, T_i)$\tcp*[f]{prefer the earliest waypoint}}{
      $t_k \gets (t_0^i + k + \Delta)\,\Delta T$; $\vect{w} \gets \vect{c}(\pi_i[k])$\;
      \lIf{$t_k - t < \dt$ \KwOr $\norm{\vect{w} - \pos_i(t)}/(t_k - t) > v_{\max}$}{\KwContinue\tcp*[f]{too fast}\label{alg:ch24-reconnect:fast}}
      \lIf{\HybridSegment{$\pos_i(t)$, $t$, $\vect{w}$, $t_k$, prediction, teammates}}{\KwRet $(k, \Delta)$\label{alg:ch24-reconnect:ok}}
    }
  }
  \KwRet \emph{failure}\label{alg:ch24-reconnect:fail}\;
}
\Fn{\HybridSegment{$\pos$, $t_a$, $\vect{w}$, $t_b$, prediction, teammates}}{
  \ForEach{sample time $t' = t_a + (t_b - t_a)\,j/n$, $j = 0, \dots, n$, with $n = \lceil (t_b - t_a)/\dt \rceil$}{
    $\vect{q} \gets \pos + (\vect{w} - \pos)(t' - t_a)/(t_b - t_a)$\;
    \lIf{clearance of $\vect{q}$ to the obstacles $< r_A$}{\KwRet \KwFalse}
    \lIf{$\norm{\vect{q} - \hat{\pos}_{T+h}} - \kappa_p\sigma(h) < R_{\mathrm{safe}}$ for $h = (t' - t)/\dt$}{\KwRet \KwFalse\label{alg:ch24-reconnect:pred}}
    \lIf{$\norm{\vect{q} - \tilde{\pos}_j(t' - t)} < R_{\mathrm{mate}}$ for some teammate $a_j$}{\KwRet \KwFalse\label{alg:ch24-reconnect:mate}}
  }
  \KwRet \KwTrue\;
}
\end{algorithm}

The search order of \cref{alg:ch24-reconnect} encodes the preferences:
delay zero first, because a reconnection with $\Delta = 0$ returns the drone
to its own reservation and nothing else in the world model changes; among
equal delays the earliest waypoint first, because it minimises the time
spent off the plan. The look-ahead $K$ bounds the work to
$(\Delta_{\max} + 1)(K + 1)$ segment tests, each of $\bigO{n\,k}$ distance
computations, and it also bounds the detour: a reconnection never aims more
than $K$ steps ahead, so the segment is at most $(K + \Delta_{\max})\ell$ long.

Two consequences matter for the rest of the design. A reconnection with
$\Delta = 0$ needs no conflict re-check: the plan in force is unchanged,
and condition (iv) has verified the segment against the teammates. A
reconnection with $\Delta > 0$ shifts the tail of $\pi_i$ by $\Delta$ steps
(line~\ref{alg:ch24-loop:shift} of \cref{alg:ch24-loop}), and the shifted
reservation may now collide with a teammate that was timed to cross behind
the drone; the executive must therefore run the re-check of
\cref{sec:ch24-replan} after every delayed reconnection. In the worked
scenario this is exactly what happens to drone~C at $t = 4.9$\,s: its
reconnection with $\Delta = 1$ succeeds, and the re-check finds the conflict
$\langle B, C, (6,2), 7 \rangle$ that the shift created.

When is a reconnection ``too costly'' or ``impossible''? Impossible when
every candidate fails line~\ref{alg:ch24-reconnect:fast} or the segment
test, as for drone~A at $t = 9.1$\,s in the scenario, where the intruder
flies along A's own lane and every segment back to it crosses the predicted
tube. Too costly when only delays beyond $\Delta_{\max}$ would work, or, in a
design that keeps a formation, when the segment would stretch a link beyond
$R_{\mathrm{comm}}$ or leave the formation slot for longer than the mission
tolerates (\cref{sec:ch24-constraints}). All of these end on
line~\ref{alg:ch24-reconnect:fail}, and \cref{alg:ch24-loop} escalates to
Replanning.

% ---------------------------------------------------------------------
\section{Replanning and the conflict re-check}
\label{sec:ch24-replan}

\paragraph{The replanning trigger.}
The executive replans in four situations, all visible in
\cref{tab:ch24-transitions}: the reconnection search failed; the drift
exceeded $d_{\max}$ while Avoiding, so that the plan has lost contact with
the drone; the drone has been Avoiding for longer than $t_{\mathrm{avoid}}$,
which is how a deadlock of the local layer (\cref{ch:ch13}) is detected; or
the local program was infeasible, meaning the half-planes left no velocity
in the disc and the dense fallback of \cref{ch:ch13} is choosing the least
bad one. A fifth trigger comes from the constraints of
\cref{sec:ch24-constraints}. None of these fires because the plan is
\emph{suboptimal}; a drone that is safe and connected keeps its plan even
if a shorter one has become available, because every replan costs a re-check
and risks a cascade (second pitfall of \cref{sec:ch24-implementation}).

\paragraph{What the replanner searches.}
The replanning layer searches the time-expanded grid of \cref{ch:ch04} from
the space-time start $(c_0, t_{\mathrm{start}})$ of \cref{sec:ch24-world}
to the drone's goal, with three kinds of obstacles:
\begin{enumerate}
  \item the \textbf{reservation table} $R$ of the teammates' \emph{current}
  paths, in absolute grid time: every (cell, time) pair, every reversed move
  (against swaps) and every parked goal from its arrival time on, exactly as
  in the low level of \cbs and in prioritized planning (\cref{ch:ch09,ch:ch08});
  \item the teammates' current cells, blocked at $t_{\mathrm{start}}$,
  because a teammate that is itself Avoiding is not where its plan says;
  \item the \textbf{blocked layers}\index{blocked layer} $\mathcal{B}_{t'}$
  of the predicted intruder for $t_{\mathrm{start}} \le t' \le t + T_{\mathrm{pred}}$,
  the cells within $R_{\mathrm{safe}} + \kappa_p\sigma(h) + \ell/2$ of the
  predicted mean at the matching step $h$; beyond the prediction horizon the
  intruder is unknown and nothing is blocked.
\end{enumerate}
A layer may be \emph{blocked} or merely \emph{expensive}. Blocking gives the
strongest statement (the path is clear of the $p$-disc at every step) but
can make the problem infeasible in a corridor; a high traversal cost keeps
the search complete and lets the drone cross the tube when there is no
alternative. The code of this chapter blocks; \cref{exr:ch24-states} asks
for the soft version. The search itself is \dstarlite (\cref{ch:ch05}) or
space-time \astar (\cref{ch:ch04}). \dstarlite is the natural choice when the
layers change every cycle and the drone keeps moving: the changed cells are
edge-cost changes, the moving start is what the key modifier $k_m$ absorbs,
and the previous search tree is reused. The simulation of this chapter
calls a fresh space-time \astar for every replan, because its instances have
a few hundred states and a fresh search costs a millisecond; the interface
(current cell, start time, reservation, blocked layers) is the same, so the
two are interchangeable.

\paragraph{Why a re-check is needed.}
The replanned path respects the reservation table, so why check it again?
Because the table is not the world. First, it holds \emph{plans}, and a
teammate that is Avoiding is off its plan; its reservation protects the
cells it will occupy once it is back, not the ones it flies through now.
Second, two drones may replan in the same cycle, each against the other's
\emph{old} path, and the two new paths can collide although each respected
what it saw. Third, a delayed reconnection shifts a path without any search
at all. The re-check therefore runs after every change of a path in force:
the conflict detector of \cref{ch:ch07} over the padded paths of all drones
from the current grid time on. If it finds a conflict, the executive does
not restart the global planner. It runs a \textbf{local
\cbs}\index{partial replanning}\index{CBS!local (partial replanning)}, also
called partial replanning, on the two agents in conflict, from
$t_{\mathrm{start}}$, with the paths of all other drones reserved; a
conflict between the pair and a third drone is impossible by construction,
and the pair's own conflicts are resolved by the constraint tree of
\cref{ch:ch09}. \Cref{alg:ch24-replan} states the whole procedure.

\begin{algorithm}[htb]
\caption{Replanning from the current state and the conflict re-check}
\label{alg:ch24-replan}
\KwIn{drone $a_i$ at $\pos_i(t)$, goal $\gamma_i$, time $t$, prediction, the paths $(\pi_j, t_0^j)$ of all drones}
\KwOut{success (paths updated) or \emph{failure}}
\Fn{\HybridReplan{$a_i$, $t$, prediction}}{
  $c_0 \gets$ the cell containing $\pos_i(t)$, or its nearest free neighbour; $t_{\mathrm{start}} \gets \lceil t/\Delta T + \tfrac12 \rceil$\label{alg:ch24-replan:start}\;
  $R \gets$ reservation of $(\pi_j, t_0^j)$ for all $j \ne i$, plus the cell of every $\pos_j(t)$ at $t_{\mathrm{start}}$\label{alg:ch24-replan:res}\;
  \ForEach{grid time $t' \ge t_{\mathrm{start}}$ with $t' \Delta T - t \le T_{\mathrm{pred}}$, step $h = (t'\Delta T - t)/\dt$}{
    add to $R$ the layer $\mathcal{B}_{t'} = \set{v \in V : \norm{\vect{c}(v) - \hat{\pos}_{T+h}} < R_{\mathrm{safe}} + \kappa_p\sigma(h) + \ell/2}$\label{alg:ch24-replan:layers}\;
  }
  $\pi' \gets$ \HybridSearch{$G$, $c_0$, $\gamma_i$, $t_{\mathrm{start}}$, $R$}\tcp*{\dstarlite or space-time \astar with the goal-stay test}\label{alg:ch24-replan:search}
  \lIf{$\pi' =$ failure}{\KwRet \emph{failure}}
  $\pi_i \gets (c_0, \pi'[0], \pi'[1], \dots)$; $t_0^i \gets t_{\mathrm{start}} - 1$\label{alg:ch24-replan:install}\;
  $\mathit{conflict} \gets$ \HybridDetect{$(\pi_j, t_0^j)_{j}$, from $\lfloor t/\Delta T \rfloor$}\tcp*{detector of \cref{ch:ch07}}\label{alg:ch24-replan:detect}
  \If{$\mathit{conflict} = \langle a_i, a_j, \dots \rangle$}{
    $R' \gets$ reservation of every drone except $a_i, a_j$\;
    $(\pi_i, \pi_j) \gets$ \HybridCbs{starts: cells of $a_i, a_j$ at $t_{\mathrm{start}}$; goals $\gamma_i, \gamma_j$; $t_{\mathrm{start}}$; $R'$}\label{alg:ch24-replan:cbs}\;
    \lIf{failure}{escalate: \cbs on a larger subset, or on all drones from their current cells\label{alg:ch24-replan:escalate}}
    $t_0^i \gets t_0^j \gets t_{\mathrm{start}}$; $a_j \gets$ Reconnecting towards $\vect{c}(\pi_j[0])$ at $t_{\mathrm{start}}$\;
  }
  \KwRet success\;
}
\end{algorithm}

\begin{pitfall}[Reservations without their time offsets]
A path that starts at $t_0 = 10$ occupies its first cell at grid time $10$,
not at time $0$. An executive that stores paths without their start times,
or that pads a replanned path backwards to time zero when it checks for
conflicts, sees conflicts that never happen (a drone that ``occupies''
its current cell during the whole past) and misses those that do (a
reservation shifted by the age of the plan). The same error appears when
grid time and clock time drift apart, for instance after a delayed
reconnection. Keep one clock, store $(\pi_i, t_0^i)$ together, convert
with \cref{eq:ch24-reference} only, and run the detector from the current
grid time, never from zero.
\end{pitfall}

% ---------------------------------------------------------------------
\section{Formation and communication constraints during avoidance and replanning}
\label{sec:ch24-constraints}

A swarm that must keep a formation or a radio link has two more
constraints, and the architecture handles them in three places.

\paragraph{In the preferred velocity.}
In the Nominal and Reconnecting states the executive adds to
\cref{eq:ch24-vpref} the displacement-based formation correction of
\cref{ch:ch23}: for a follower $a_i$ with leader $a_\ell$ and desired
displacement $\vect{d}_{\ell i}$,
\begin{equation}
  \vel^{\mathrm{pref}}_i \;\gets\; \vel^{\mathrm{pref}}_i + k_F\bigl(\pos_\ell + \vect{d}_{\ell i} - \pos_i\bigr),
  \label{eq:ch24-formation}
\end{equation}
which is one term of the consensus protocol on the shifted positions
$\pos_i - \vect{o}_i$; the general form sums over the formation
neighbours. The communication range enters as a soft barrier that acts
before the link breaks: if the nearest teammate $a_j$ is farther than an
activation length $\ell_{\mathrm{act}} < R_{\mathrm{comm}}$,
\begin{equation}
  \vel^{\mathrm{pref}}_i \;\gets\; \vel^{\mathrm{pref}}_i + \frac{\pos_j - \pos_i}{\norm{\pos_j - \pos_i}}\bigl(\norm{\pos_j - \pos_i} - \ell_{\mathrm{act}}\bigr),
  \label{eq:ch24-comm}
\end{equation}
and the result is clipped to $v_{\max}$. Both terms are inputs to the local
layer, not overrides of it: \orca may still change the velocity, and when it
does, safety wins.

\paragraph{What is relaxed while Avoiding, and restored while Reconnecting.}
Separation from the intruder has priority over the shape of the swarm. In
the Avoiding state the executive drops the formation term
\cref{eq:ch24-formation}: keeping the slot would pull the drone back into the
conflict it is leaving, and it would also drag the whole formation after a
single avoiding drone through the followers' corrections. The range term
\cref{eq:ch24-comm} stays, because a broken link is a mission failure rather
than a shape error, and it is soft enough not to fight the half-planes. The
teammates that are not avoiding keep their formation terms and therefore
move to accommodate the one that is, as \cref{ch:ch23} describes. When the
drone enters Reconnecting the formation gain returns, in the code at once,
in a smoother design ramped over a second, and the formation error $e_F$ of
\cref{ch:ch23} decays with the consensus rate. With the MPC tracker of
\cref{ch:ch21} the same policy is expressed in the constraint rows: the
formation slot and the range are soft constraints with slacks, the slack
weight of the formation is set to zero while Avoiding and restored while
Reconnecting, and the range constraint is kept hard whenever the program
stays feasible. The recovery time of $e_F$ after the intruder has passed is
one of the metrics of \cref{ch:ch25}.

\paragraph{In the replanning trigger and the re-check.}
A reconnection or a replanned path is a single-drone decision, and it can
be admissible in every geometric sense and still break the swarm: the
straight segment of \cref{alg:ch24-reconnect} may take the drone beyond
$R_{\mathrm{comm}}$ of every teammate, and a replanned path may keep it
there for ten steps. The executive therefore adds a constraint test to both:
over the horizon of the prediction, for every grid time $t'$, the padded
paths must keep at least one link of every drone shorter than
$R_{\mathrm{comm}}$ (or, for a formation, keep the formation graph
connected), and the formation error of the planned positions must return
below its tolerance $e_{\max}$ within the recovery time. A candidate that
fails the test counts as ``violates the constraints'' in the training
plan's rule (iv): the reconnection is rejected and the drone replans; a
replanned path that fails it is rejected and the executive escalates to
line~\ref{alg:ch24-replan:escalate} of \cref{alg:ch24-replan} with the
formation planned as a group, for instance as the virtual leader of
\cref{ch:ch23} whose footprint is the formation's. In the worked scenario
the range is monitored only: the longest link between a drone and its
nearest teammate over the run is $7.33$\,m, below $\ell_{\mathrm{act}} = 7.5$\,m,
so \cref{eq:ch24-comm} never acts and the communication graph stays
connected throughout.

% ---------------------------------------------------------------------
\section{Prediction-aware constraints with uncertainty}
\label{sec:ch24-uncertainty}

Every decision above uses the prediction through one number, the inflation
$\kappa_p\,\sigma(h)$. This section says where it comes from and what it
guarantees.

\paragraph{From a covariance to a radius.}
The prediction layer delivers, at step $h$, a Gaussian
$\Normal(\hat{\pos}_{T+h}, \mat{\Sigma}_h)$ for the intruder's position
(\cref{ch:ch18,ch:ch20}). The set of positions within Mahalanobis distance
$\kappa$ of the mean is the ellipse $\mathcal{E}_h(\kappa)$ of \cref{ch:ch02},
and in two dimensions the probability mass inside it is
$1 - e^{-\kappa^2/2}$, because the squared Mahalanobis distance of a
two-dimensional Gaussian is exponentially distributed with mean $2$. Hence
the ellipse of confidence $p$ has
\begin{equation}
  \kappa_p = \sqrt{-2\ln(1 - p)},\qquad \kappa_{0.95} = 2.45,\quad \kappa_{0.99} = 3.03.
  \label{eq:ch24-kappa}
\end{equation}
The disc of radius $\kappa_p\sigma(h)$, with $\sigma(h) = \sqrt{\lambda_{\max}(\mat{\Sigma}_h)}$,
contains this ellipse (\cref{ch:ch02}: the semi-axes are
$\kappa_p\sqrt{\lambda_i}$), so the intruder's centre lies inside it with
probability at least $p$, and the \textbf{inflated safety
radius}\index{inflated radius}\index{prediction!uncertainty inflation}
\begin{equation}
  d^{\mathrm{safe}}_h = R_{\mathrm{safe}} + \kappa_p\,\sigma(h) = r_A + r_B + \Delta r + \kappa_p\sqrt{\lambda_{\max}(\mat{\Sigma}_h)}
  \label{eq:ch24-inflated}
\end{equation}
is the distance from the predicted mean at which the drone's disc is clear
of the intruder's disc with probability at least $p$. This is the formula
of \cref{ch:ch20}, where the factor is written $k_p$. The chapter's code
uses $p = 0.95$.

\paragraph{The chance-constraint reading.}
Requiring $\norm{\pos_A - \hat{\pos}_{T+h}} \ge d^{\mathrm{safe}}_h$ is a
conservative version of the \textbf{chance constraint}\index{chance
constraint} $\Prob(\norm{\pos_A - \pos_B} \ge R_{\mathrm{safe}}) \ge p$ of
\textcite{zhu2019chance}, which \cref{ch:ch21} imposes in its linear form
with the standard deviation taken along the normal of the half-plane and the
one-dimensional factor $\kappa_\delta = 1.64$ for $\delta = 0.05$. The disc
form is what the trigger and the replanner need, because neither has a
normal direction to project on; it costs a factor $2.45/1.64$ in radius and
guarantees the same probability. In both forms the guarantee is per step
and per obstacle: over $H$ steps and $n$ intruders the risk adds up to at
most $H n (1 - p)$, which is why $p$ is chosen close to one and why the
safety horizon is short.

\paragraph{Three consumers.}
The same sequence of discs is read three ways. The trigger and the local
layer use one disc: \cref{alg:ch24-trigger} inflates at every look-ahead,
and once the trigger has fired the \orca half-plane of the intruder is built
with the radius $r_B + \Delta r + \kappa_p\sigma(h_c)$ at the step $h_c$ of
the predicted collision, with full responsibility because the intruder does
not reciprocate (\cref{ch:ch13}). The replanner uses one layer of cells per
grid time, \cref{alg:ch24-replan} line~\ref{alg:ch24-replan:layers}; in the
worked scenario the replan of drone~A at $t = 9.1$\,s blocks $72$ cells over
five layers, nine of them in the first. The MPC tracker of \cref{ch:ch21}
uses one half-plane per step with the inflated radius. A predictor that
reports a wrong covariance corrupts all three: the LSTM of \cref{ch:ch20}
covered $54\,\%$ instead of $95\,\%$ at $2.4$\,s before calibration, and an
over-confident predictor makes the trigger fire late, the layers too thin
and the MPC margin too small. Calibrate the ellipses per horizon on
validation data before they enter this chapter.

\paragraph{Particles instead of ellipses.}
When the prediction is a set of $N$ weighted trajectories, from the
particle filter of \cref{ch:ch19} or from a multimodal predictor of
\cref{ch:ch20}, the disc is replaced by a count. The trigger fires at
look-ahead $s$ when the total weight of the particles $\pos^{(n)}_h$ with
$\norm{\tilde{\pos}_i(s) - \pos^{(n)}_h} < R_{\mathrm{safe}}$ exceeds
$1 - p$; a cell enters the layer $\mathcal{B}_{t'}$ when the weight of the
particles within $R_{\mathrm{safe}} + \ell/2$ of its centre exceeds the same
threshold; and the \orca disc is built around the mode that threatens the
drone, with a radius that covers its particles. This handles an intruder
that may pass on either side of a building without inflating one ellipse
over the whole block, at a cost linear in $N$ per test.
\Cref{exr:ch24-particles} works this out.

\paragraph{How fast the inflation grows.}
For the constant-velocity Kalman model with white-noise acceleration
(\cref{ch:ch18}), the position variance of an extrapolation grows with the
cube of the look-ahead, so $\sigma(h)$ grows like $s^{3/2}$. In the
scenario the tracker reports $\sigma = 0.06, 0.14, 0.24, 0.35$\,m at
$s = 0, 1, 2, 3$\,s, hence an inflation of $0.86$\,m at the horizon, larger
than $R_{\mathrm{safe}}$ itself. The uncertainty, not the geometry, decides
the trigger at the horizon, and a longer horizon with a constant-velocity
prediction would mostly add false alarms; a learned predictor with a
calibrated, slower-growing covariance is what a longer horizon buys, which
is one of the research directions of \cref{sec:ch24-research}.
